% !PRINTED 586
\weberpage{586}
Entwickelt man die linke Seite nach Potenzen von $w$, so sind
die Koeffizienten rational durch $\lambda$ und $\sqrt{\Delta}$ ausdrückbar, gehören
also dem Körper
\[
\mathfrak{L} = \mathfrak{R}(\lambda, \sqrt{\Delta})
\]
an. Entwickelt man die rechte Seite in gleicher Weise und
ordnet nach den Potenzen von $w$, so tritt $A_\nu$ zuerst in dem
Koeffizienten von $w^{\nu-1}$ auf, und zwar mit dem Faktor 1 behaftet.
Danach läßt sich $A_\nu$ bestimmen, wenn die früheren Koeffizienten
schon bestimmt sind, und es folgt, daß alle diese Koeffizienten
dem Körper $\mathfrak{L}$ angehören.

Da man nachträglich wieder Zähler und Nenner des Bruches
$A/D$ durch rationale Rechnung von gemeinschaftlichen Faktoren
befreien kann, so ergibt sich der Satz:

2. \gesperrt{Wenn in dem Ausdruck (8) Zähler und Nenner
$A(\boldsymbol{s})$, $D(\boldsymbol{s})$ von gemeinschaftlichen Faktoren befreit
sind, so gehören diese Funktionen dem Körper $\mathfrak{L}$ an.}
% [source sic: the theorem cites expression (8), although A and D were introduced in (10)]

Um die Funktion $A(x)$ darzustellen, nehmen wir den Koeffizienten
 der höchsten Potenz von $x$ gleich 1 an und untersuchen,
für welche Werte von $u$ die Funktion $\boldsymbol{s}(\mu u)$ verschwindet. Wir
bezeichnen einen Wert $u$, für den $\boldsymbol{s}(\mu u)$ verschwindet, mit $2\varrho/\mu$
und erhalten dann als Bedingung des Verschwindens nach (5)
eine der beiden Kongruenzen
\[
2\varrho \equiv 0, 1 \quad (\mathrm{mod}\ 2, \omega).
\]
Es ist also $\varrho$ eine Zahl des Körpers $\Omega$, die zwar gebrochen
sein, aber keine anderen Nenner als einen Teiler von $2a$ haben
kann. Andererseits kann man $\varrho$ um ein Vielfaches von $\mu$ verändern,
ohne daß der Wert $\boldsymbol{s}(2\varrho/\mu)$ geändert wird, und da es
nur auf den letzteren ankommt, und man der Kongruenz:
\[
\varrho \equiv \frac{\xi}{2a} \quad (\mathrm{mod}\ \mu),
\]
wenn $\xi$ eine ganze Zahl in $\Omega$ ist, durch eine ganze Zahl $\varrho$ genügen
kann, so können wir $\varrho$ als \gesperrt{ganze Zahl} annehmen, wenn wir \gesperrt{voraussetzen,
daß der erste Koeffizient $a$ in (1) ungerade und
relativ prim zu $\mu$ sei.} Diese Annahme halten wir von jetzt an
fest.

Lassen wir also $\varrho$ ein vollständiges Restsystem nach dem
Modul $\mu$ mit Ausschluß der Null durchlaufen, so erhalten wir
alle Werte von $\boldsymbol{s}$, für die $A(\boldsymbol{s})$ verschwindet, in der Form:
\weq{13}{\boldsymbol{s}\left(\frac{2\varrho}{\mu}\right).}
